% =====================================================================
% Appendix D (app:baselines): baselines, wrappers, operating points, prompts
% (one-column appendix). Float: Table A4 (tables/tab_baselines.tex).
% Implemented: the offline systems with oracle input (offline baseline script
% of the cst-bench package; prompt, max_new_tokens 160 and batch 16 quoted
% from it). Planned: the streaming variant with oracle input, both wrappers,
% the streaming baselines, the decoupled twin and the operating points (each
% marked with \draftnote). Systems with oracle input are not called bounds.
% =====================================================================
\section{Baselines and Wrappers}\label{app:baselines}

\Cref{tab:baselines} lists the translation systems of \cref{sec:exp}. All of them are scored by \scorer from committed, time-stamped pieces (\cref{app:scorer}); this appendix describes how each system turns a session into pieces, how its latency settings are chosen, and which prompts the LLM-based systems use. We use public checkpoints without fine-tuning on \task data; the only tuning data is the \syn calibration split, and nothing is tuned on \taxi. Seed LiveInterpret~2.0 \citep{cheng2025seed}, a closed system, covers only Chinese--English, which \bench lacks, so we discuss it but do not run it.
\draftnote{implemented: the offline systems with oracle input, whose \taxi scores are preliminary; planned: the streaming variant with oracle input, both wrappers, the streaming baselines, the decoupled twin and the operating points}

% =====================================================================
% Table A4 (tab:baselines): baseline configurations. Included from App. D.
% Known cells only: checkpoints and parameter counts given by the model names
% or the model cards, offline settings, directions, wrappers, and the
% licences already verified (Qwen3.8-27B, Qwen3-ASR, Nemotron 3.5 ASR,
% Hibiki-Zero). Everything else, including both values of every latency
% knob, is \tbd. System names follow Table 3 (tab:main) and Table 6
% (tab:ablation). Systems with oracle input are not called bounds.
% 600M (model card) and 638M (measured) count the whole Nemotron 3.5 ASR
% model, RNN-T decoder included, not its encoder alone.
% =====================================================================
\begin{table}[tp]
  \caption{Baseline configurations, named as in \cref{tab:main,tab:ablation} (LLM: Qwen3.8-27B). Params: recognition plus translation models (B: billions; \model: total; its encoder comes from Nemotron~3.5 ASR Streaming 0.6B, which has 600M parameters per its model card and 638M as measured, RNN-T decoder included). Directions: \pair{En}{De} on \taxi and, for systems with a Korean entry, on \syn; \pair{Ko}{En} on \syn only. Policy: when text is committed; LA-2: LocalAgreement-2. Latency knob: the setting swept on the \syn calibration split; its two values (low and high latency) are chosen by the rule of \cref{app:baselines} and are pending. Wrappers: O oracle segmentation, R realistic (VAD and language-ID routing), --~none (the system handles the whole mix). Not listed: the diarization cascades of \cref{tab:aux}, which do not translate (\cref{app:baselines}). $^\ast$Optional; InfiniSST only if its weights are public; Hibiki-Zero is scored on its text output. Blank cells are pending.}
  \label{tab:baselines}
  \centering
  \fontsize{8}{9.5}\selectfont % 8pt (ICML \footnotesize is 9pt): fits the width
  \setlength{\tabcolsep}{2.5pt}
  \begin{tabular}{@{}llllllll@{}}
    \toprule
    System & Checkpoint\,/\,version & Params & Directions & Policy & \makecell[l]{Latency knob\\(two points)} & Wrappers & Licence \\
    \midrule
    \multicolumn{8}{@{}l}{\textit{Offline, oracle input: gold turns, each translated after it ends}} \\
    Gold$\to$LLM
      & Qwen3.8-27B & 27B
      & \makecell[l]{\pair{En}{De}\\\pair{Ko}{En}} & turn end & -- & O & Apache-2.0 \\
    ASR$\to$LLM
      & \makecell[l]{Qwen/Qwen3-ASR-1.7B;\\Qwen3.8-27B} & 1.7B\,+\,27B
      & \makecell[l]{\pair{En}{De}\\\pair{Ko}{En}} & turn end & -- & O & Apache-2.0 \\
    Whisper ST
      & openai/whisper-large-v3 & \tbd
      & \makecell[l]{\dir{De}{En}\\\dir{Ko}{En}} & turn end & -- & O & \tbd \\
    \midrule
    \multicolumn{8}{@{}l}{\textit{Streaming, oracle input: gold transcripts with forced-aligned word times; non-causal unit choice}} \\
    Gold$\to$LLM (MU2)
      & Qwen3.8-27B & 27B
      & \pair{En}{De} & MU2 units & \makecell[l]{one point\\(chrF $\geq$ 50)} & O & Apache-2.0 \\
    \midrule
    \multicolumn{8}{@{}l}{\textit{Pipelines built from single-direction systems}} \\
    \makecell[l]{Consecutive\\cascade}
      & \makecell[l]{Silero VAD; LID;\\Qwen3-ASR-1.7B;\\Qwen3.8-27B} & 1.7B\,+\,27B
      & \makecell[l]{\pair{En}{De}\\\pair{Ko}{En}} & VAD endpoint & \makecell[l]{endpoint\\silence} & R & \tbd \\
    SeamlessStreaming
      & \tbd & \tbd
      & \makecell[l]{\pair{En}{De}\\\pair{Ko}{En}} & EMMA & \makecell[l]{decision\\threshold} & O, R & \tbd \\
    2$\times$SeamlessStreaming
      & as above & \tbd
      & \pair{En}{De} & EMMA & as above & -- & \tbd \\
    \makecell[l]{Streaming\\cascade}
      & \makecell[l]{nvidia/nemotron-3.5-\\asr-streaming-0.6b;\\Qwen3.8-27B} & 0.6B\,+\,27B
      & \makecell[l]{\pair{En}{De}\\\pair{Ko}{En}} & LA-2 & \makecell[l]{ASR chunk,\\update step} & O, R & \makecell[l]{OpenMDW-1.1;\\Apache-2.0} \\
    M4T v2 + AlignAtt
      & SeamlessM4T v2 & \tbd
      & \pair{En}{De} & AlignAtt & frames $f$ & O, R & \tbd \\
    \makecell[l]{Canary-1B-v2\\+ AlignAtt$^\ast$}
      & Canary-1B-v2 & 1B
      & \pair{En}{De} & AlignAtt & frames $f$ & O, R & \tbd \\
    StreamSpeech
      & \tbd & \tbd
      & \dir{De}{En} & learned & chunk size & O, R & \tbd \\
    Hibiki-Zero$^\ast$
      & \makecell[l]{hibiki-zero-3b-\\pytorch-bf16} & 3B
      & \dir{De}{En} & learned & \tbd & O, R & \makecell[l]{CC BY-NC-SA\\4.0} \\
    InfiniSST$^\ast$
      & \tbd & \tbd
      & \dir{En}{De} & learned (LLM) & \tbd & O, R & \tbd \\
    \midrule
    \multicolumn{8}{@{}l}{\textit{Controls in the ablations (\cref{tab:ablation})}} \\
    Decoupled twin
      & \makecell[l]{\model M1; separate\\1.7B-thinker pass} & \tbd
      & \pair{En}{De} & \makecell[l]{MU2 units,\\delay $\delta_t$} & $\delta_t$ & -- & \makecell[l]{OpenMDW-1.1;\\Apache-2.0} \\
    \makecell[l]{Same-backbone cascade\\(M1$\to$LLM)}
      & \makecell[l]{\model M1;\\Qwen3.8-27B} & 2,336M\,+\,27B
      & \pair{En}{De} & LA-2 & update step & -- & \makecell[l]{OpenMDW-1.1;\\Apache-2.0} \\
    \makecell[l]{Size-matched cascade\\(ASR$\to$1.7B LLM)}
      & \makecell[l]{Qwen3-ASR-1.7B;\\LLM of $\approx$1.7B} & 1.7B\,+\,$\approx$1.7B
      & \pair{En}{De} & LA-2 & update step & R & \tbd \\
    \midrule
    \multicolumn{8}{@{}l}{\textit{Ours}} \\
    \model
      & \makecell[l]{Nemotron 3.5 encoder;\\Qwen3-ASR-1.7B\\thinker} & 2,336M
      & \makecell[l]{\pair{En}{De}\\\pair{Ko}{En}} & \makecell[l]{MU2 units,\\delay $\delta_t$} & $\delta_t$ & -- & \makecell[l]{OpenMDW-1.1;\\Apache-2.0} \\
    \bottomrule
  \end{tabular}
\end{table}


\paragraph{Offline systems with oracle input.}
Each gold turn is cut from \texttt{mix.wav} at its reference interval, 0.1\,s margins included, and translated after it ends, one turn per prompt and without dialogue history. Qwen3-ASR-1.7B \citep{shi2026qwen3asr} transcribes the turn with its language given, and Qwen3.8-27B translates this transcript, or the gold transcript, from the user message
\begin{quote}\small\ttfamily\raggedright
Translate this \{src\_lang\} utterance from a phone conversation into \{tgt\_lang\}. Output only the \{tgt\_lang\} translation.
\end{quote}
followed by a blank line and the turn's text; the placeholders are English names of the languages, such as \texttt{German}. Decoding is greedy with at most 160 new tokens and the model's reasoning mode switched off, and the first line of the answer is kept. Whisper-large-v3 \citep{radford2023robust} translates German turns, and on \syn Korean turns, into English with its translation task and the source language given; it cannot translate into German. All models process batches of 16 turns. Every piece carries its turn ID and is stamped at the turn end $e_n$, so on the ideal clock \eo is 0 and \laal equals $L_n$ (\cref{app:scorer}). The preliminary \taxi scores of these systems were computed in the scorer's oracle mode, by turn ID and without resegmentation, and will be re-scored with resegmentation under the final protocol. We do not call these systems bounds: they translate without dialogue history, and on \deen Whisper ST already scores above Gold$\to$LLM in BLEU (\cref{tab:main}). The computation-aware clock adds the per-turn compute, measured as the time of a batch divided by its size and therefore approximate; \cref{tab:ca-latency} re-measures it without batching. The WER of the ASR transcripts is computed after the scorer's normalisation.

\paragraph{Streaming variant with oracle input.}
To show what the commit policy costs and saves when recognition and segmentation are perfect, Gold$\to$LLM (MU2) reveals gold transcripts with forced-aligned word times (Qwen3-ForcedAligner; \citealp{shi2026qwen3asr}) word by word to Qwen3.8-27B, which commits translations at verified meaning units (MU2; \cref{sec:model}), accepting a unit if its chrF against the aligned part of the full translation is at least 50. A unit is stamped when its last source word ends, and the rest of a turn is flushed at the turn end. As units are chosen with the whole turn in view, this system is non-causal; it runs only under the oracle wrapper and so never becomes the reference pipeline $P$ of \cref{sec:exp-setup}.
\draftnote{planned; the German minimum unit is set on held-out data before the run}

\paragraph{Oracle wrapper.}
A single-direction system runs as two instances, one per target language. The wrapper cuts each gold turn $[b_n,e_n]$ from the mono mix, so that under L1 the cut also contains the overlapping speech of the other speaker, and streams it in real time to the instance that translates into the other speaker's language. Instances are reset at every turn and receive no dialogue history. A piece emitted after the instance has consumed $u$ seconds of the turn is stamped $b_n+u$ on the ideal clock, plus compute time on the computation-aware clock, and carries the turn ID. The main scores are computed with the turn IDs removed, so that the scorer resegments oracle-wrapper output like any other; scoring with them (oracle mode, \cref{app:scorer}) is a diagnostic of how much resegmentation alone moves the scores. Cutting the turn from its speaker's channel of \texttt{2ch.wav} instead, in the channel order recorded in the timeline, gives the oracle-separation variant of \cref{tab:factorial}.
\draftnote{planned}

\paragraph{Realistic wrapper.}
The realistic wrapper receives only the mix and the language pair. Silero VAD \citep{silero2024silero} segments the stream online: a segment opens at a detected speech onset and closes after a minimum silence. A spoken-language identifier restricted to the session's two languages labels each segment from its onset and updates the label as audio arrives. The segment is streamed to the instance that translates into the other language and is re-routed to the other instance if the label flips; text already committed stays. Pieces are stamped from the segment onset as in the oracle wrapper. As identifier we use the automatic language detection of Nemotron~3.5 ASR \citep{nvidia2026nemotron35asr}, or the language identification of Qwen3-ASR \citep{shi2026qwen3asr} if the former proves less accurate on the calibration split. The VAD's speech-probability threshold, minimum silence and speech padding and the identifier's settings are fixed once on that split, for all systems. As the split is clean 16\,kHz TTS, these settings may suit telephone-band \taxi less well (\cref{app:syn}). A short minimum silence splits turns at pauses, whereas a long one merges turns across the short L0 natural gaps (median drawn gap 0.223\,s on \taxi, to which the 0.1\,s turn margins add up to 0.2\,s of silence; \cref{sec:bench-timing}); under L1, the VAD hears one continuous segment across the speaker change, so the second turn inherits the first turn's route unless the label flips. Sweeps of the minimum silence on \taxi show how sensitive the realistic wrapper is to this setting; they serve analysis only, not tuning. The factorial of \cref{tab:factorial} combines components of the two wrappers: gold turns with the identifier, and VAD segments with the gold language of the turn that overlaps each segment most.
\draftnote{planned; the identifier is chosen on the calibration split}

\paragraph{Streaming systems.}
SeamlessStreaming \citep{seamless2023seamless} translates speech into a requested language with Efficient Monotonic Multihead Attention (EMMA), a variant of monotonic multihead attention \citep{ma2020monotonic}, and runs behind both wrappers. Two SeamlessStreaming instances on the whole mix, one per target language and without routing, form the naive bidirectional baseline: each instance also receives speech that is already in its target language, and what it emits for that speech counts as wrong-direction output (copy-through or wrong language; \cref{app:scorer}). The streaming cascade passes the growing transcript of Nemotron~3.5 ASR \citep{nvidia2026nemotron35asr}, a cache-aware streaming RNN-T, to Qwen3.8-27B, which re-translates it after every update, with the committed translation as a forced prefix and the instruction of the offline prompt, and commits the longest common prefix of its last two translations (LocalAgreement-2; \citealp{liu2020lowlatency,polak2022cuni,machacek2023turning}). It follows recent cascades that translate a growing transcript with an LLM \citep{fuxa2026alignatt4llm}, but its policy needs no access to attention weights. SeamlessM4T~v2 \citep{seamless2023seamless,seamless2025joint} and, optionally, Canary-1B-v2 \citep{sekoyan2025canary} are offline speech translation models that we run simultaneously in simulstream \citep{gaido2025simulstream} with AlignAtt \citep{papi2023alignatt}, which emits a token only if the audio frame it attends to most is not among the last $f$ frames. StreamSpeech \citep{zhang2024streamspeech}, which learns translation and policy jointly, covers \deen; InfiniSST \citep{ouyang2025infinisst}, an LLM-based system for unbounded speech, covers \ende and is included only if its weights are public. Hibiki-Zero \citep{labiausse2026simultaneous}, a simultaneous speech-to-speech and speech-to-text translation model trained without word-level aligned data, translates German among other languages into English and has public weights (CC BY-NC-SA 4.0); it is an optional \deen baseline, scored on its text output only.
\draftnote{planned; checkpoints and parameter counts left blank in \cref{tab:baselines} are filled when the runs are set up}

\paragraph{Cascades.}
The consecutive cascade is the realistic counterpart of the offline cascade: it waits for the VAD endpoint of each segment, identifies the segment's language and then runs Qwen3-ASR-1.7B and Qwen3.8-27B with the offline prompt, so its \eo includes the endpoint silence and the compute time. In \cref{tab:ablation}, the same-backbone cascade feeds the transcripts of M1, \model's stage-2 checkpoint (\cref{sec:model}), to the LocalAgreement-2 translation stage, and the size-matched cascade pairs Qwen3-ASR-1.7B with an LLM of about 1.7B parameters, the size of \model's language model. The diarization cascades of \cref{tab:aux} combine pyannote \citep{bredin2023pyannote} with offline Qwen3-ASR-1.7B, and Streaming Sortformer \citep{medennikov2025streaming} with streaming Nemotron~3.5 ASR.
\draftnote{planned}

\paragraph{Decoupled twin.}
The decoupled twin of \cref{tab:ablation} tests the coupling of turn-taking and translation at \model's scale and commit policy. It starts from the same M1 checkpoint and 1.7B thinker as \model: M1 writes the speaker-tagged transcript, whose speaker tags fix turn boundaries and directions as hard decisions, and a separate translation pass, trained on the same data, then translates each fixed segment with the same MU2 commit policy and $\delta_t$. Translation therefore cannot revise segmentation or routing, while scale, data, simulated-session exposure and commit policy match \model; the same-backbone cascade shares the M1 exposure but not the translator's scale or policy.
\draftnote{planned; the twin needs its own Stage-3 run}

\paragraph{Operating points.}
Each system with a latency knob runs at two operating points, chosen once on the \syn calibration split (about 5\% of its sessions) by a rule that is fixed before the test runs and is the same for all systems. We sweep the knob under the oracle wrapper, so that one setting serves both wrappers, and for each of two budgets $\Lambda_{\mathrm{low}}<\Lambda_{\mathrm{high}}$ on the mean \laal (\tbdtext) take the setting with the highest COMET (chrF for Korean targets) that stays within the budget; a system that cannot meet a budget runs at its fastest setting, which we mark. The consecutive cascade, which exists only in realistic form, sweeps the minimum silence of its own VAD instead, and \model's two operating points are values of $\delta_t$ chosen by the same rule. The offline systems and their streaming variant with oracle input have a single operating point.
\draftnote{planned; until \syn is built, systems run at their documented default settings and are re-run once the calibration split exists; decision needed: whether the low- or the high-budget setting is the main operating point}

\paragraph{Release.}
We release the wrapper code with the scorer, including segmentation, language identification, routing and the settings of every operating point, so that any single-direction system can be evaluated under the same protocol. For \taxi, we release per-turn scores keyed by turn ID but no system outputs (\cref{app:repro}).
